\section{Conclusion and Future Work}
\label{sec:conclusion}

\subsection{Summary of Technical Contribution}
This documentation has presented the complete architectural, deployment, and operational specification for \textbf{Research Rover}. By uniting multi-source literature retrieval (PubMed and OpenAlex), neural semantic re-ranking (MedCPT), dense vector indexing (Sentence Transformers and FAISS), and retrieval-augmented conversational generation (Google Gemini), Research Rover offers a comprehensive environment for data-driven scientific inquiry.

The development and integration of the \textbf{Universal Dataset Ingestion and Quality Analysis System} represents a transformative enhancement to the platform's reliability. By introducing:
\begin{enumerate}
    \item A resilient \textbf{20-field canonical schema} with a four-tier header matching cascade,
    \item Multi-attribute \textbf{quality scoring} (0--100) with duplicate and invalid value deductions,
    \item Deterministic \textbf{seven-feature readiness evaluation} guarding downstream analytics, and
    \item Transparent \textbf{frontend telemetry integration} via modals and banners,
\end{enumerate}
Research Rover effectively eliminates brittle column dependencies, shields bibliometric algorithms from corrupt inputs, and provides researchers with complete clarity regarding dataset hygiene before initiating computational workflows.

\subsection{Directions for Future Extension}
Subsequent engineering iterations may explore the following extensions:
\begin{itemize}
    \item \textbf{Interactive Header Remapping:} Expanding the frontend Dataset Quality Modal to allow researchers to manually correct or assign unmapped columns and persist those overrides into the sidecar.
    \item \textbf{Expanded File Formats:} Extending the ingestion pipeline beyond delimited CSV files to natively parse BibTeX (\texttt{.bib}), Research Information Systems (\texttt{.ris}), and EndNote XML formats.
    \item \textbf{Automated Index Synchronization:} Incorporating cryptographic hash validation between stored FAISS vector indexes and parent CSV datasets to automatically flag stale embeddings when records are appended.
    \item \textbf{Distributed Vector Stores:} Transitioning FAISS CPU indexes to distributed vector databases (e.g., Milvus or Qdrant) for multi-million document enterprise scaling.
\end{itemize}
